% ECONOMICS_PROBLEM: {"id": "I11-551", "title": "Provider payment and quality", "jel_code": "I11", "source_id": "OP-0053"}
% !TeX program = xelatex
\documentclass[11pt,a4paper]{article}
\usepackage[margin=23mm]{geometry}
\usepackage{fontspec}
\setmainfont{Latin Modern Roman}
\usepackage{amsmath,amssymb,mathtools}
\usepackage{xcolor,enumitem,booktabs,longtable,array,xurl,hyperref}
\definecolor{muted}{HTML}{53636C}
\hypersetup{colorlinks=true,urlcolor=blue,linkcolor=blue}
\setlength{\parindent}{0pt}
\setlength{\parskip}{5pt}
\newcommand{\R}{\mathbb R}
\newcommand{\E}{\mathbb E}
\newcommand{\N}{\mathbb N}
\newcommand{\ind}{\mathbf 1}
\newcommand{\Var}{\operatorname{Var}}
\newcommand{\Cov}{\operatorname{Cov}}
\newcommand{\supp}{\operatorname{supp}}
\DeclareMathOperator*{\argmin}{arg\,min}
\DeclareMathOperator*{\argmax}{arg\,max}
\newcommand{\entrymeta}[3]{{\sffamily\small\color{muted}\textbf{\texttt{#1}}\quad|\quad #2\par #3\par}}
\newcommand{\Problemref}[1]{\texttt{#1}}
\newcommand{\JELref}[2]{\texttt{#2} (JEL \texttt{#1})}
\begin{document}
\section*{Provider payment and quality}
\textbf{Problem ID:} \texttt{I11-551}\quad \textbf{JEL:} \texttt{I11}\par
% BEGIN ECONOMICS STATEMENT
\label{problem:OP-0053}
\label{atlas:prob:H3}
\noindent\textbf{Original atlas identifier:} H3; entry 53. \textbf{Field:} Health economics. \textbf{Original tractability label:} Medium.

\noindent\textbf{Classification:} Parameterized theoretical/empirical research agenda; not a certified unsolved theorem.

\subsection*{Definitions and assumptions}
Let provider $h$ choose treatment intensity $a$, coding $k$, quality effort $q$, referrals, and patient acceptance after observing information $\mathcal I_h$. A payment contract $p$ maps verifiable diagnoses, services, and outcomes into reimbursement. A structural model $m$ specifies patient risk, provider preferences, clinical production, costs, reporting behavior, and selected patient/provider equilibrium. Define $H_i(p;m)$ as the health outcome for each patient in a fixed eligible population, including patients denied or diverted from care; $v_i\geq0$ converts health to a monetary welfare measure. Let $C(p;m)$ denote real treatment, quality, administrative, and coding-resource costs, and $A(p;m)$ access-related travel/time losses. Budget expenditure $B(p;m)$ includes transfers to providers but is not automatically a social resource cost.

\subsection*{Formal research question}
\emph{Editorial formulation: the mathematical target below is not a theorem quoted from the references.}

For contracts $p\in\mathcal P$ satisfying the payer budget, provider participation, verifiability, and an explicit minimum-access/quality constraint, characterize
\[
 \sup_{p\in\mathcal P}\left\{\mathbb E_m\sum_i w_i v_iH_i(p;m)-C(p;m)-A(p;m)\right\},
\]
where $w_i\geq0$ are fixed distributional weights. If provider rents or fiscal distortions enter welfare, add them with declared weights and count transfers consistently. Determine which fee and bundled-payment variation identifies true treatment substitution separately from coding, risk selection, and referrals. For mixed prospective/service-based contracts, report the identified set of clinical and budget effects and $\varepsilon$-optimal contracts conditional on provider objectives. Lower measured volume alone is not the formal efficiency criterion.

\subsection*{Interpretation and scope}
Unnecessary volume is defined by marginal expected health value relative to real opportunity cost, not by a comparison with an administratively chosen average. Stinting means a reduction in beneficial care or access relative to that criterion and must be measured with outcomes beyond coded claims. Payments can affect who is observed, so evaluating treated patients alone omits selection and delayed treatment. Risk adjustment may protect high-risk access yet induce diagnosis inflation; those mechanisms need separate measures. Provider altruism is a structural parameter rather than a universal empirical constant. Quasi-experimental fees identify responses in a particular reimbursement regime; extrapolating to capitation requires a model of effort and selection. The extension targets payment rules that jointly control real costs and clinical losses under feasible monitoring.

\subsection*{Source audit and status}
The original three works are identifiable and provide existing supply-side models and payment evidence. [S2] is correctly titled “under Prospective Payment,” not a generic Medicare-system title. They do not prove universal dominance of any payment form, nor establish that the editorial constrained optimization is a currently open theorem.

\subsection*{References}
\begin{enumerate}[label={[S\arabic*]},leftmargin=*,itemsep=0.6em]
\item Randall P. Ellis and Thomas G. McGuire (1986). Provider Behavior under Prospective Reimbursement: Cost Sharing and Supply. Journal of Health Economics 5(2), 129--151. DOI: 10.1016/0167-6296(86)90002-0.
\url{https://www.sciencedirect.com/science/article/pii/0167629686900020}
\newline\emph{Locator and support limits:} Abstract and model: prospective payment can induce insufficient services; no universally dominant payment rule.
\item David M. Cutler (1995). The Incidence of Adverse Medical Outcomes under Prospective Payment. Econometrica 63(1), 29--50.
\url{https://cutler.scholars.harvard.edu/publications/incidence-adverse-medical-outcomes-under-prospective-payment}
\newline\emph{Locator and support limits:} Author publication record and article: prospective hospital payment and adverse outcomes. NBER antecedent is WP4300, March 1993.
\item Jeffrey Clemens and Joshua D. Gottlieb (2014). Do Physicians’ Financial Incentives Affect Medical Treatment and Patient Health? American Economic Review 104(4), 1320--1349. DOI: 10.1257/aer.104.4.1320.
\url{https://www.aeaweb.org/articles?id=10.1257/aer.104.4.1320}
\newline\emph{Locator and support limits:} Article and publisher metadata: Medicare fee changes, treatment, and outcomes; not evaluation of every mixed-payment contract.
\end{enumerate}

\subsection*{Common conventions: Source mathematical conventions}

Unless an entry states otherwise, agent and firm sets are finite. State and action spaces are measurable, strategies and outcome maps are measurable, probability laws are specified, and expectations exist for the targets under discussion. Infinite-horizon discount factors lie in $(0,1)$ and discounted objectives are finite. Norms, tolerances and complexity input sizes must be specified for computational comparisons. Constants or primitives that appear locally are inputs, not empirical facts asserted by this catalogue.

\textbf{Identification.} A statistical model $m\in\mathcal M$ implies an observable law $P_m$ and target $\theta(m)$. At law $P$, its identified set is
\[
\Theta(P;\mathcal M)=\{\theta(m):m\in\mathcal M,\ P_m=P\}.
\]
Point identification holds when this set is a singleton, or equivalently when $P_{m_1}=P_{m_2}$ implies $\theta(m_1)=\theta(m_2)$ for all compatible models. Sharp bounds mean bounds attained, or approached when the boundary is not attained, by compatible models. Writing this set is a definition, not a solution: the substantive task is to establish informative restrictions and derive or estimate the set. An unrestricted-confounding model may give only trivial bounds.

\textbf{Inference.} Identification, estimability and valid uncertainty quantification are separate questions. Fix $\alpha\in(0,1)$. A confidence procedure $C_n$ over a declared law class $\mathcal P$ has asymptotic uniform coverage at level $1-\alpha$ when
\[
\liminf_{n\to\infty}\inf_{P\in\mathcal P}\Pr_P\{\theta(P)\in C_n\}\geq1-\alpha.
\]
For set-identified targets the coverage event must be defined separately, for example coverage of the full identified set. If an entry uses identification-indexed models instead of uniquely indexed laws, the same coverage convention is applied over those models. Pointwise limits do not establish uniform validity, and asymptotic validity does not guarantee finite-sample precision. Sample sizes, dependence structures and weak-identification sequences are part of each application.

\textbf{Counterfactuals.} $Y_i(a)$ is a potential outcome under a specified intervention, including its timing, implementation and versions. It is not $Y_i$ conditional on voluntarily choosing $a$. A full intervention vector permits interference; shorthand scalar treatment notation does not silently impose no interference. Assignment, selection, attrition, anticipation and measurement must be specified. A claim about an instrument's local effect is not automatically a population policy effect.

\textbf{Economic equilibrium.} A counterfactual model includes best responses, constraints, feasibility, market clearing, state transitions and expectations consistent with the stated information. When several equilibria exist, either a declared selector is an input or the target is set-valued. Comparisons across policy regimes must use the stated selection convention. Reduced-form relationships are not presumed invariant after an equilibrium-changing intervention.

\textbf{Optimization and welfare.} Optimization is over the stated feasible set. If attainment is not established, $\sup$ or $\inf$ and $\varepsilon$-optimal choices replace an asserted maximizer or minimizer. For example, with value $V=\sup_{a\in\mathcal A}W(a)<\infty$, an $\varepsilon$-optimal set is $\{a\in\mathcal A:W(a)\geq V-\varepsilon\}$ for $\varepsilon>0$. Benefits, costs, risk and health or environmental outcomes can be aggregated only using declared units and conversion weights. Interpersonal utility weights, the treatment of fiscal transfers and foreign profits, discounting, and the behavioral welfare criterion are explicit inputs. A comparative-static derivative additionally requires existence and appropriate differentiability of the selected counterfactual.

\textbf{Sources.} Local labels [S1], [S2], and so forth restart in every question. Locators identify the relevant abstract, model, section or result at the level actually checked; a bibliographic or abstract-level verification is not represented as inspection of an entire proof. Working papers are distinguished from later publications and changing versions. Source summaries are paraphrased. All URL checks and editorial formulations refer to the audit date stated above.
% END ECONOMICS STATEMENT
\end{document}
